\input{../tex-inputs/latex-leseansicht-vorspann}

\section[Rudolf Schwarzkopf an Arthur Schnitzler, 25. 9. 1893]{L04276 Rudolf Schwarzkopf an Arthur Schnitzler, 25. 9. 1893}
\nopagebreak\mylabel{L04276v}
\rehead{ }\normalsize\beginnumbering\briefempfaengerindex{Schnitzler, Arthur@\textsc{Schnitzler, Arthur}!zzzSchwarzkopf, Rudolf@\emph{von Rudolf Schwarzkopf}!1893-09-252@{25. 9. 1893}|(be}
\toendnotes[C]{\smallbreak\pagebreak[2]}
\correspDesc{Versand  durch Rudolf Schwarzkopf am 25. 9. 1893 in Meran
\newline{}Erhalt  durch Arthur Schnitzler am 26. 9. 1893 in Wien}\toendnotes[C]{\smallbreak}
\Standort{CUL, Schnitzler, B 96a.}
\physDesc{Postkarte, 657 Zeichen
\newline{}Handschrift: schwarze Tinte, deutsche Kurrent
\newline{}Versand: 1) Stempel: »\nobreak{}\oindex{Meran@\textbf{Meran}, \emph{Hauptstadt}|pwk}M\textcolor{gray}{era}{[}n{]}, {[}25{]}/9 9\textcolor{gray}{3}, 2–N\nobreak{}«.   2) Stempel: »\nobreak{}\oindex{I., Innere Stadt@\textbf{I., Innere Stadt}, \emph{Verwaltungsgebiet}|pwk}Wien 1/1, 26. 9. 93, 11½–1N, Bestellt\nobreak{}«. 
\newline{}Schnitzler: mit Bleistift bei der Datumsangabe die Jahreszahl ergänzt: »93« }\pstart{}{\pb}Herrn D\textsuperscript{or} Arthur \textsc{Schnitzler}\pend{}\pstart{}\textsc{Wien\oindex{Wien@\textbf{Wien}, \emph{Verwaltungsgebiet}|pw}}\pend{}\pstart{}\textsc{I Grillparzerſtr 7\oindex{Wien@\textbf{Wien}!I., Innere Stadt@\textbf{I., Innere Stadt}!Wohnung und Ordination Arthur Schnitzler Grillparzerstraße 7/3. Stock@\textbf{Wohnung und Ordination Arthur Schnitzler Grillparzerstraße 7/3. Stock}, \emph{Ordination}|pw}}\pend{}{\bigskip}\vspace{1em}
\pstart{}{\pb}Sehr geehrter Herr
                  Doctor!\pend\vspace{0.5em}
\pstart
           Wie Sie mittlerweile wahrſcheinlich schon von \textsc{Gustav\pwindex{Schwarzkopf, Gustav 7.\,11.\,1853 Wien – 13.\,11.\,1939 ebd.@\textsc{Schwarzkopf, Gustav} (7.\,11.\,1853 Wien – 13.\,11.\,1939 ebd.), \emph{Schriftsteller}|pw}} erfahren haben dürften, bin ich wolbehalten hie angelangt. Mein Befinden iſt,
                  \textcolor{gray}{bi}s auf die leidigen Kehlkopfſchmerzen, die nicht aufhören
               wollen, ein befriedingendes. Die Geſellſchaft hier läßt{ }ſehr viel zu wünſchen übrig;
               ganz im Gegenſatz zu unſerem Tiſch in der \textsc{Brühl\oindex{Brühl@\textbf{Brühl}, \emph{Tal}|pw}}, bin ich an unſerm \textsc{Table-d’hote}  der jüngſte. Es
               giebt alſo nicht den Stoff für das kleinſte Novellchen.\pend
           
\pstart
           Ich will hoffen daſs ich bald recht Günſtiges über \substVorne{}\textsuperscript{d}\substDazwischen{}m\substHinten{}ein Befinden berichten kann und verbleibe einſtweilen\pend
           
\pstart
           Ihr ganz ergeben{\\[\baselineskip]}\spacefill\mbox{RudolfSchwarzkopf}\pend
           \leftskip=0em{}
\pstart
           \noindent{}\textsc{Meran\oindex{Meran@\textbf{Meran}, \emph{Hauptstadt}|pw}}\pend
           
\pstart
           \raggedleft{}\textsc{Hotel Paſſerhof\oindex{Hotel {\kaufmannsund} Pension Passerhof@\textbf{Hotel {\kaufmannsund} Pension Passerhof}, \emph{Hotel}|pw}}{ }25/9\pend
           \selectlanguage{ngerman}\endnumbering\briefempfaengerindex{Schnitzler, Arthur@\textsc{Schnitzler, Arthur}!zzzSchwarzkopf, Rudolf@\emph{von Rudolf Schwarzkopf}!1893-09-252@{25. 9. 1893}|)be}\mylabel{L04276h}
\begin{anhang}
\end{anhang}\newcommand{\dateiname}{L04276}\newcommand{\titel}{Rudolf Schwarzkopf an Arthur Schnitzler, 25. 9. 1893}\newcommand{\editorInnen}{Herausgegeben von Jahnke, SelmaMüller, Martin Anton}\input{../tex-inputs/latex-leseansicht-abspann}
